\ifdefined\gkver\else\def\gkver{student}\fi
\documentclass[\gkver]{gaokaozhenti}
\begin{document}

\chapter{2016年天津理综}
%% paperType: 理综物理部分

\begin{enumerate}
\item
%% number: 1
%% paperName: 2016年天津理综第 1 题 6 分
%% typeId: 1
%% type: 单选题
%% chapter: 光学
%% point: 光的电磁说
%% method: 无
%% score: 6
%% degree: 850
%% duplicateId: 0
%% body:
我国成功研发的反隐身先进米波雷达堪称隐身飞机的克星，它标志着我国雷达研究又创新的里程碑，米波雷达发射无线电波的波长在 $1\sim10\Um$ 范围内，则对该无线电波的判断正确的是\xzanswer{C}

\onepicture[h=3.4cm]{figs/01.jpg}

\fourchoices[answer=C]
{米波的频率比厘米波频率高}
{和机械波一样须靠介质传播}
{同光波一样会发生反射现象}
{不可能产生干涉和衍射现象}

%% answer: C
%% memo:
\memoanswer{
根据 $f = \frac{v}{\lambda}$ 可知，波长越长的波频率越低，故米波的频率比厘米波的频率低，选项 A 错误；无线电波不需要介质传播，选项 B 错误；同光波一样会发生反射，选项 C 正确；干涉和衍射是波特有的现象，选项 D 错误。
}

\item
%% number: 2
%% paperName: 2016年天津理综第 2 题 6 分
%% typeId: 1
%% type: 单选题
%% chapter: 光学
%% point: 光的干涉
%% method: 无
%% score: 6
%% degree: 650
%% duplicateId: 0
%% body:
右图是 a、b 两光分别经过同一双缝干涉装置后在屏上形成的干涉图样，则\xzanswer{D}

\onepicture[h=3.2cm]{figs/02.png}

\fourchoices[answer=D]
{在同种均匀介质中，a 光的传播速度比 b 光的大}
{从同种介质射入真空发生全反射时 a 光临界角大}
{照射在同一金属板上发生光电效应时，a 光的饱和电流大}
{若两光均由氢原子能级跃迁产生，产生 a 光的能级能量差大}

%% answer: D
%% memo:
\memoanswer{
由图可知 a 光的干涉条纹间距小于 b 光的，根据 $\Delta x = \frac{L}{d} \lambda$ 可知，a 的波长小于 b 的波长，a 光的频率大于 b 光的频率，a 光的折射率大于 b 光的折射率，则根据 $n = \frac{c}{v}$ 可知，在同种介质中传播时 a 光的传播速度较小，选项 A 错误；根据 $\sin C = \frac{1}{n}$ 可知，从同种介质中射入真空，a 光发生全反射的临界角小，选项 B 错误；发生光电效应时饱和光电流与入射光的强度有关，故无法比较饱和光电流的大小，选项 C 错误；a 光的频率较高，若两光均由氢原子能级跃迁产生，则产生 a 光的能级差大，选项 D 正确。
}

\item
%% number: 3
%% paperName: 2016年天津理综第 3 题 6 分
%% typeId: 1
%% type: 单选题
%% chapter: 曲线运动
%% point: 人造地球卫星
%% method: 无
%% score: 6
%% degree: 550
%% duplicateId: 0
%% body:
我国即将发射“天宫二号”空间实验室，之后发生“神舟十一号”飞船与“天宫二号”对接。假设“天宫二号”与“神舟十一号”都围绕地球做匀速圆周运动，为了实现飞船与空间实验室的对接，下列措施可行的是\xzanswer{C}

\onepicture[h=2.8cm]{figs/03.png}

\fourchoices[answer=C]
{使飞船与空间实验室在同一轨道上运行，然后飞船加速追上空间实验室实现对接}
{使飞船与空间实验室在同一轨道上运行，然后空间实验室减速等待飞船实现对接}
{飞船先在比空间实验室半径小的轨道上加速，加速后飞船逐渐靠近空间实验室，两者速度接近时实现对接}
{飞船先在比空间实验室半径小的轨道上减速，减速后飞船逐渐靠近空间实验室，两者速度接近时实现对接}

%% answer: C
%% memo:
\memoanswer{
若使飞船与空间实验室在同一轨道上运行，飞船加速会进入较高的轨道，空间实验室减速会进入较低的轨道，都不能实现对接，选项 AB 错误；要想实现对接，可使飞船在比空间实验室半径较小的轨道上加速，然后飞船将进入较高的空间实验室轨道，逐渐靠近空间站后，两者速度接近时实现对接，选项 C 正确，选项 D 错误。
}

\item
%% number: 4
%% paperName: 2016年天津理综第 4 题 6 分
%% typeId: 1
%% type: 单选题
%% chapter: 电路
%% point: 电容
%% method: 无
%% score: 6
%% degree: 550
%% duplicateId: 0
%% body:
如图所示，平行板电容器带有等量异种电荷，与静电计相连，静电计金属外壳和电容器下极板都接地。在两极板间有一个固定在 P 点的点电荷，以 $E$ 表示两板间的电场强度，$E_{p}$ 表示点电荷在 P 点的电势能，$\theta$ 表示静电计指针的偏角。若保持下极板不动，将上极板向下移动一小段距离至图中虚线位置，则\xzanswer{D}

\onepicture[h=3.0cm]{figs/04.png}

\fourchoices[answer=D]
{$\theta$ 增大，$E$ 增大}
{$\theta$ 增大，$E_{p}$ 不变}
{$\theta$ 减小，$E_{p}$ 增大}
{$\theta$ 减小，$E$ 不变}

%% answer: D
%% memo:
\memoanswer{
若保持下极板不动，将上极板向下移动一小段距离，则根据 $C = \frac{\varepsilon S}{4\pi kd}$ 可知，$C$ 变大，$Q$ 一定，则根据 $Q = CU$ 可知，$U$ 减小，则静电计指针偏角 $\theta$ 减小；根据 $E = \frac{U}{d}$，$Q = CU$，$C = \frac{\varepsilon S}{4\pi kd}$，联立可得 $E = \frac{4\pi kQ}{\varepsilon S}$，可知 $Q$ 一定时，$E$ 不变；根据 $U_1 = Ed_1$ 可知 P 点离下极板的距离不变，$E$ 不变，则 P 点与下极板的电势差不变，P 点的电势不变，则 $E_p$ 不变；故选项 ABC 错误，D 正确。
}

\item
%% number: 5
%% paperName: 2016年天津理综第 5 题 6 分
%% typeId: 1
%% type: 单选题
%% chapter: 交流电
%% point: 变压器
%% method: 无
%% score: 6
%% degree: 450
%% duplicateId: 0
%% body:
如图所示，理想变压器原线圈接在交流电源上，图中各电表均为理想电表。下列说法正确的是\xzanswer{B}

\onepicture[h=3.4cm]{figs/05.png}

\fourchoices[answer=B]
{当滑动变阻器的滑动触头 P 向上滑动时，$R_{1}$ 消耗的功率变大}
{当滑动变阻器的滑动触头 P 向上滑动时，电压表 V 示数变大}
{当滑动变阻器的滑动触头 P 向上滑动时，电流表 $A_{1}$ 示数变大}
{若闭合开关 S，则电流表 $A_{1}$ 示数变大、$A_{2}$ 示数变大}

%% answer: B
%% memo:
\memoanswer{
当滑动变阻器的滑动触头 P 向上滑动时，$R$ 的阻值变大，副线圈中电流变小，原线圈中电流也变小，电流表 $A_{1}$ 示数变小；$R_{1}$ 消耗的功率及两端电压均变小，副线圈总电压不变，则电压表的示数变大，选项 AC 错误，B 正确。若闭合开关 S，则副线圈总电阻变小，电流变大，原线圈中电流变大，电流表 $A_{1}$ 示数变大；$R_{1}$ 两端电压变大，$R_{2}$ 两端电压变小，电流表 $A_{2}$ 示数变小，选项 D 错误。
}

\item
%% number: 6
%% paperName: 2016年天津理综第 6 题 6 分
%% typeId: 2
%% type: 多选题
%% chapter: 其他
%% point: 物理学史
%% method: 无
%% score: 6
%% degree: 650
%% duplicateId: 0
%% body:
物理学家通过对实验的深入观察和研究，获得正确的科学认知，推动物理学的发展。下列说法符合事实的是\xzanswer{AC}

\fourchoices[answer=AC]
{赫兹通过一系列实验，证实了麦克斯韦关于光的电磁理论}
{查德威克用 $\alpha$ 粒子轰击 $\ce{^{14}_{7}N}$ 获得反冲核 $\ce{^{17}_{8}O}$，发现了中子}
{贝克勒尔发现的天然放射性现象，说明原子核有复杂结构}
{卢瑟福通过对阴极射线的研究，提出了原子核式结构模型}

%% answer: AC
%% memo:
\memoanswer{
麦克斯韦预言了电磁波的存在，赫兹通过实验证实了麦克斯韦的电磁理论，选项 A 正确；卢瑟福用 $\alpha$ 粒子轰击 $\ce{^{14}_{7}N}$，获得反冲核 $\ce{^{17}_{8}O}$，发现了质子，选项 B 错误；贝克勒尔发现的天然放射性现象，说明原子核具有复杂结构，选项 C 正确；卢瑟福通过对 $\alpha$ 粒子散射实验的研究，提出了原子的核式结构模型，选项 D 错误。
}

\item
%% number: 7
%% paperName: 2016年天津理综第 7 题 6 分
%% typeId: 2
%% type: 多选题
%% chapter: 机械波
%% point: 波的图象
%% method: 无
%% score: 6
%% degree: 450
%% duplicateId: 0
%% body:
在均匀介质中坐标原点 O 处有一波源做简谐运动，其表达式为 $y = 5\sin\left(\frac{\pi}{2}t\right)$，它在介质中形成的简谐横波沿 $x$ 轴正方向传播，某时刻波刚好传播到 $x = 12\Um$ 处，波形图象如图所示，则\xzanswer{AB}

\onepicture[h=3.4cm]{figs/07.png}

\fourchoices[answer=AB]
{此后再经 $6\Us$ 该波传播到 $x = 24\Um$ 处}
{M 点在此后第 $3\Us$ 末的振动方向沿 $y$ 轴正方向}
{波源开始振动时的运动方向沿 $y$ 轴负方向}
{此后 M 点第一次到达 $y = -3\Um$ 处所需时间是 $2\Us$}

%% answer: AB
%% memo:
\memoanswer{
波的周期 $T=4\Us$，波长 $\lambda=8\Um$，波速 $v=\frac{\lambda}{T}=2\Ums$，则再经过 $6\Us$，波传播的距离为 $x=vt=12\Um$，故该波传到 $x=24\Um$ 处，选项 A 正确；M 点在此时振动方向向下，则第 $3\Us$ 末，即经过了 $0.75T$，该点的振动方向沿 $y$ 轴正方向，选项 B 正确；因波传到 $x=12\Um$ 处时，质点向 $y$ 轴正方向振动，故波源开始振动时的运动方向沿 $y$ 轴正方向，选项 C 错误；因为 $\frac{3}{5}<\frac{\sqrt{2}}{2}=\sin 45^\circ$，故 M 点第一次到达 $y=-3\Um$ 位置时，振动的时间小于 $\frac{T}{4}=1\Us$，选项 D 错误。
}

\item
%% number: 8
%% paperName: 2016年天津理综第 8 题 6 分
%% typeId: 2
%% type: 多选题
%% chapter: 运动和力
%% point: 连接体
%% method: 无
%% score: 6
%% degree: 450
%% duplicateId: 0
%% body:
我国高铁技术处于世界领先水平，和谐号动车组是由动车和拖车编组而成，提供动力的车厢叫动车，不提供动力的车厢叫拖车。假设动车组各车厢质量均相等，动车的额定功率都相同，动车组在水平直轨道上运行过程中阻力与车重成正比。某列动车组由 8 节车厢组成，其中第 1、5 节车厢为动车，其余为拖车，则该动车组\xzanswer{BD}

\onepicture[h=3.0cm]{figs/08.jpg}

\fourchoices[answer=BD]
{启动时乘客受到车厢作用力的方向与车运动的方向相反}
{做匀加速运动时，第 5、6 节与第 6、7 节车厢间的作用力之比为 $3:2$}
{进站时从关闭发动机到停下来滑行的距离与关闭发动机时的速度成正比}
{与改为 4 节动车带 4 节拖车的动车组最大速度之比为 $1:2$}

%% answer: BD
%% memo:
\memoanswer{
列车启动时，乘客随车厢加速运动，加速度方向与车的运动方向相同，故乘客受到车厢的作用力方向与车运动方向相同，选项 A 错误。动车组运动的加速度 $a = \frac{2F - 8kmg}{8m} = \frac{F}{4m} - kg$，对第 6、7、8 节车厢整体：$F_{56} = 3ma + 3kmg = 0.75F$；对第 7、8 节车厢整体：$F_{67} = 2ma + 2kmg = 0.5F$；故第 5、6 节车厢与第 6、7 节车厢间的作用力之比为 $3:2$，选项 B 正确。根据动能定理 $\frac{1}{2}Mv^{2} = kMgs$，解得 $s = \frac{v^{2}}{2kg}$，可知进站时从关闭发动机到停下来滑行的距离与关闭发动机时速度的平方成正比，选项 C 错误；8 节车厢有 2 节动车时的最大速度为 $v_{m1} = \frac{2P}{8kmg}$；8 节车厢有 4 节动车时的最大速度为 $v_{m2} = \frac{4P}{8kmg}$，则 $\frac{v_{m1}}{v_{m2}} = \frac{1}{2}$，选项 D 正确。
}

\item
%% number: 9
%% paperName: 2016年天津理综第 9 题 3 分
%% typeId: 4
%% type: 实验题
%% chapter: 电路
%% point: 伏安法测电阻
%% method: 无
%% score: 3
%% degree: 450
%% duplicateId: 0
%% body:
\begin{enumerate}
	\item
	如图所示，方盒 A 静止在光滑的水平面上，盒内有一个小滑块 B，盒的质量是滑块质量的 2 倍，滑块与盒内水平面间的动摩擦因数为 $\mu$。若滑块以速度 $v$ 开始向左运动，与盒的左右壁发生无机械能损失的碰撞，滑块在盒中来回运动多次，最终相对盒静止，则此时盒的速度大小为\tkanswer[1.4cm]{$\frac{v}{3}$}，滑块相对于盒运动的路程为\tkanswer[1.8cm]{$\frac{v^{2}}{3\mu g}$}。
	\onepicture[w=7.5cm, align=right]{figs/09.png}
	\item
	某同学利用如图 \ref{2016天津理综09a} 所示的装置研究小车的匀变速直线运动。
	\begin{steps}
		\item
		实验中必要的措施是\tkanswer[1.2cm]{AB}。

		（A）细线必须与长木板平行

		（B）先接通电源再释放小车

		（C）小车的质量远大于钩码的质量

		（D）平衡小车与长木板间的摩擦力
		\item
		他实验时将打点计时器接到频率为 $50\UHz$ 的交流电源上，得到一条纸带，打出的部分计数点如图 \ref{2016天津理综09b} 所示（每相邻两个计数点间还有 4 个点，图中未画出）。$s_{1}=3.59\Ucm$，$s_{2}=4.41\Ucm$，$s_{3}=5.19\Ucm$，$s_{4}=5.97\Ucm$，$s_{5}=6.78\Ucm$，$s_{6}=7.64\Ucm$。则小车的加速度 $a=$\tkanswer[1.4cm]{0.80}$\Umsq$（要求充分利用测量的数据），打点计时器在打 B 点时小车的速度 $v_{B}=$\tkanswer[1.4cm]{0.40}$\Ums$。（结果均保留两位有效数字）
	\end{steps}
	\twopicture[labelA=2016天津理综09a, labelB=2016天津理综09b, align=right, hA=3.0cm, wB=12cm]
	{figs/09a.png}
	{figs/09b.png}
	\item
	某同学想要描绘标有“$3.8\UV$，$0.3\UA$”字样小灯泡 L 的伏安特性曲线，要求测量数据尽量精确，绘制曲线完整，可供该同学选用的器材除了开关、导线外，还有：

	电压表 $V_{1}$（量程 $0\sim3\UV$，内阻等于 $3\UkO$）

	电压表 $V_{2}$（量程 $0\sim15\UV$，内阻等于 $15\UkO$）

	电流表 $A_{1}$（量程 $0\sim200\UmA$，内阻等于 $10\UO$）

	电流表 $A_{2}$（量程 $0\sim3\UA$，内阻等于 $0.1\UO$）

	滑动变阻器 $R_{1}$（$0\sim10\UO$，额定电流 $2\UA$）

	滑动变阻器 $R_{2}$（$0\sim1\UkO$，额定电流 $0.5\UA$）

	定值电阻 $R_{3}$（阻值等于 $1\UO$）

	定值电阻 $R_{4}$（阻值等于 $10\UO$）

	定值电阻 $R_{5}$（阻值等于 $1\UkO$）

	电源 $E$（$E=6\UV$，内阻不计）
	\begin{steps}
		\item
		请画出实验电路图，并将各元件字母代码标在该元件的符号旁。
		\item
		该同学描绘出的 $I-U$ 图象应是下图中的\tkanswer[1.2cm]{B}。

		\fourchoices[answer=B, ispicture=true, h=2.2cm]
		{figs/09c.png}
		{figs/09d.png}
		{figs/09e.png}
		{figs/09f.png}
	\end{steps}
\end{enumerate}

%% answer: 
\jdanswer*{
\begin{enumerate}
	\item
	$\frac{v}{3}$，$\frac{v^{2}}{3\mu g}$

	\item
	\begin{enumerate}
		\item
		AB

		\item
		$0.80$，$0.40$

	\end{enumerate}

	\item
	\begin{enumerate}
		\item
		如图所示（电路图略）

		\item
		B

	\end{enumerate}

\end{enumerate}
}
%% memo:
\memoanswer{
（1）设滑块质量为 $m$，则盒的质量为 $2m$；对整个过程，由动量守恒定律可得 $mv = 3mv_{\text{共}}$，解得 $v_{\text{共}} = \frac{v}{3}$，由能量守恒定律可知 $\mu mgx = \frac{1}{2}mv^{2} - \frac{1}{2} \cdot 3m \cdot \left(\frac{v}{3}\right)^{2}$，解得 $x = \frac{v^{2}}{3\mu g}$。

（2）①实验时，细线必须与长木板平行，以减小实验的误差，选项 A 正确；实验时要先接通电源再释放小车，选项 B 正确；此实验中没必要使小车的质量远大于钩码的质量，选项 C 错误；此实验中不需平衡小车与长木板间的摩擦力，选项 D 错误。

②两相邻计数点间的时间间隔 $T = 0.1\Us$；由逐差法可得 $a = \frac{s_6 + s_5 + s_4 - s_3 - s_2 - s_1}{9T^2} = 0.80\Umsq$，打点计时器在打 B 点时小车的速度 $v_B = \frac{s_1 + s_2}{2T} = 0.40\Ums$。

（3）①用电压表 $V_1$ 和 $R_2$ 串联，可改装成量程为 $U = \frac{U_{V_1}}{r_{V_1}} (r_{V_1} + R_2) = 4\UV$ 的电压表；用电流表 $A_1$ 与 $R_4$ 并联可改装为量程为 $I = I_{A_1} + \frac{I_{A_1} r_{A_1}}{R_4} = 0.4\UA$ 的电流表；待测小灯泡的阻值较小，故采用电流表外接法；为使曲线完整，滑动变阻器应采用分压接法，故选择总阻值小的 $R_1$，电路如图。

②小灯泡灯丝的电阻随温度升高而变大，故该同学描绘出的 $I-U$ 图线应该是 B。
}

\item
%% number: 10
%% paperName: 2016年天津理综第 10 题 16 分
%% typeId: 6
%% type: 计算题
%% chapter: 曲线运动
%% point: 竖直平面圆周运动
%% method: 无
%% score: 16
%% degree: 650
%% duplicateId: 0
%% body:
我国将于 2022 年举办冬奥会，跳台滑雪是其中最具观赏性的项目之一。如图所示，质量 $m = 60\Ukg$ 的运动员从长直助滑道 AB 的 A 处由静止开始以加速度 $a = 3.6\Umsq$ 匀加速滑下，到达助滑道末端 B 时速度 $v_{B} = 24\Ums$，A 与 B 的竖直高度差 $H = 48\Um$。为了改变运动员的运动方向，在助滑道与起跳台之间用一段弯曲滑道衔接，其中最低点 C 处附近是一段以 O 为圆心的圆弧。助滑道末端 B 与滑道最低点 C 的高度差 $h = 5\Um$，运动员在 B、C 间运动时阻力做功 $W = -1530\UJ$，取 $g = 10\Umsq$。
\begin{enumerate}
	\item
	求运动员在 AB 段下滑时受到阻力 $F_{f}$ 的大小；
	\item
	若运动员能够承受的最大压力为其所受重力的 6 倍，则 C 点所在圆弧的半径 $R$ 至少应为多大。
\end{enumerate}
\onepicture[h=4.0cm, align=right]{figs/10.png}

%% answer: 
\jdanswer{
\begin{enumerate}
	\item
	$144\UN$

	\item
	$12.5\Um$

\end{enumerate}
}
%% memo:
\memoanswer{
（1）运动员在 AB 上做初速度为零的匀加速运动，设 AB 的长度为 $x$，则有 $v_{B}^{2} = 2ax$ ①

由牛顿第二定律有 $\frac{H}{x}mg - F_{f} = ma$ ②

联立①②式，代入数据解得 $F_{f}=144\UN$ ③

（2）设运动员到达 C 点时的速度为 $v_{C}$，在由 B 到达 C 的过程中，由动能定理有

$mgh + W = \frac{1}{2} m v_{C}^{2} - \frac{1}{2} m v_{B}^{2}$ ④

设运动员在 C 点所受的支持力为 $F_{N}$，由牛顿第二定律有 $F_{N}-mg=m\frac{v_{C}^{2}}{R}$ ⑤

由运动员能够承受的最大压力为其所受重力的 6 倍，联立④⑤式，代入数据解得 $R = 12.5\Um$ ⑥
}

\item
%% number: 11
%% paperName: 2016年天津理综第 11 题 18 分
%% typeId: 6
%% type: 计算题
%% chapter: 磁场
%% point: 带电粒子在复合场中的运动
%% method: 无
%% score: 18
%% degree: 450
%% duplicateId: 0
%% body:
如图所示，空间中存在着水平向右的匀强电场，电场强度大小为 $E = 5\sqrt{3}\UNC$，同时存在着水平方向的匀强磁场，其方向与电场方向垂直，磁感应强度大小 $B = 0.5\UT$。有一带正电的小球，质量 $m = 1 \times 10^{-6}\Ukg$，电荷量 $q = 2 \times 10^{-6}\UC$，正以速度 $v$ 在图示的竖直面内做匀速直线运动，当经过 P 点时撤掉磁场（不考虑磁场消失引起的电磁感应现象），取 $g = 10\Umsq$。求：
\begin{enumerate}
	\item
	小球做匀速直线运动的速度 $v$ 的大小和方向；
	\item
	从撤掉磁场到小球再次穿过 P 点所在的这条电场线经历的时间 $t$。
\end{enumerate}
\onepicture[h=3.4cm, align=right]{figs/11.png}

%% answer: 
\jdanswer{
\begin{enumerate}
	\item
	$20\Ums$，与电场方向夹角为 $60^\circ$

	\item
	$3.5\Us$

\end{enumerate}
}
%% memo:
\memoanswer{
（1）小球匀速直线运动时受力如图，其所受的三个力在同一平面内，合力为零，有

$qvB = \sqrt{q^{2}E^{2} + m^{2}g^{2}}$ ①

代入数据解得 $v = 20\Ums$ ②

\onepicture[w=4.2cm]{figs/11_an.jpg}

速度 $v$ 的方向与电场 $E$ 的方向之间的夹角 $\theta$ 满足 $\tan\theta=\frac{qE}{mg}$ ③

代入数据解得 $\tan\theta=\sqrt{3}$，$\theta=60^{\circ}$ ④

（2）解法一：撤去磁场，小球在重力与电场力的合力作用下做类平抛运动，设其加速度为 $a$，有

$a=\frac{\sqrt{q^{2}E^{2}+m^{2}g^{2}}}{m}$ ⑤

设撤掉磁场后小球在初速度方向上的分位移为 $x$，有 $x = vt$ ⑥

设小球在重力与电场力的合力方向上分位移为 $y$，有 $y = \frac{1}{2} at^{2}$ ⑦

$a$ 与 $mg$ 的夹角和 $v$ 与 $E$ 的夹角相同，均为 $\theta$，又 $\tan\theta=\frac{y}{x}$ ⑧

联立④⑤⑥⑦⑧式，代入数据解得 $t=2\sqrt{3}\Us=3.5\Us$ ⑨

解法二：撤去磁场后，由于电场力垂直于竖直方向，它对竖直方向的分运动没有影响，以 P 点为坐标原点，竖直向上为正方向，小球在竖直方向上做匀减速运动，其初速度为 $v_{0}=v\sin\theta$

若使小球再次穿过 P 点所在的电场线，仅需小球的竖直方向上分位移为零，则有 $v_{0}t - \frac{1}{2}gt^{2} = 0$

联立⑤⑥式，代入数据解得 $t=2\sqrt{3}\Us=3.5\Us$ ⑦。
}

\item
%% number: 12
%% paperName: 2016年天津理综第 12 题 20 分
%% typeId: 6
%% type: 计算题
%% chapter: 电磁感应
%% point: 电磁感应中的变加速
%% method: 无
%% score: 20
%% degree: 350
%% duplicateId: 0
%% body:
电磁缓速器是应用于车辆上以提高运行安全性的辅助制动装置，其工作原理是利用电磁阻尼作用减缓车辆的速度。电磁阻尼作用可以借助如下模型讨论：如图所示，将形状相同的两根平行且足够长的铝条固定在光滑斜面上，斜面与水平方向夹角为 $\theta$。一质量为 $m$ 的条形磁铁滑入两铝条间，恰好匀速穿过，穿过时磁铁两端面与两铝条的间距始终保持恒定，其引起电磁感应的效果与磁铁不动、铝条相对磁铁运动相同。磁铁端面是边长为 $d$ 的正方形，由于磁铁距离铝条很近，磁铁端面正对两铝条区域的磁场均可视为匀强磁场，磁感应强度为 $B$，铝条的高度大于 $d$，电阻率为 $\rho$。为研究问题方便，铝条中只考虑与磁铁正对部分的电阻和磁场，其他部分电阻和磁场可忽略不计，假设磁铁进入铝条间以后，减少的机械能完全转化为铝条的内能，重力加速度为 $g$。
\begin{enumerate}
	\item
	求铝条中与磁铁正对部分的电流 $I$；
	\item
	若两铝条的宽度均为 $b$，推导磁铁匀速穿过铝条间时速度 $v$ 的表达式；
	\item
	在其他条件不变的情况下，仅将两铝条更换为宽度 $b'>b$ 的铝条，磁铁仍以速度 $v$ 进入铝条间，试简要分析说明磁铁在铝条间运动时的加速度和速度如何变化。
\end{enumerate}
\onepicture[h=6.2cm, align=right]{figs/12.png}

%% answer: 
\jdanswer{
\begin{enumerate}
	\item
	$I=\frac{mg\sin \theta }{2Bd}$

	\item
	$v=\frac{\rho mg\sin \theta }{2B^{2}d^{2}b}$

	\item
	磁铁做加速度逐渐减小的减速运动，直到 $F'=mg\sin\theta$ 时，磁铁重新达到平衡状态，将再次以较小的速度匀速下滑。

\end{enumerate}
}
%% memo:
\memoanswer{
（1）磁铁在铝条间运动时，两根铝条受到的安培力大小相等均为 $F_{\text{安}}$，有 $F_{\text{安}} = IdB$ ①

磁铁受到沿斜面向上的作用力为 $F$，其大小有 $F = 2F_{\text{安}}$ ②

磁铁匀速运动时受力平衡，则有 $F - mg \sin \theta = 0$ ③

联立①②③式可得 $I = \frac{mg \sin \theta}{2Bd}$ ④

（2）磁铁穿过铝条时，在铝条中产生的感应电动势为 $E$，有 $E = Bdv$ ⑤

铝条与磁铁正对部分的电阻为 $R$，由电阻定律有 $R = \rho \frac{d}{db}$ ⑥

由欧姆定律有 $I = \frac{E}{R}$ ⑦

联立④⑤⑥⑦式可得 $v = \frac{\rho mg \sin \theta}{2B^{2}d^{2}b}$ ⑧

（3）磁铁以速度 $v$ 进入铝条间，恰好做匀速运动时，磁铁受到沿斜面向上的作用力 $F$，联立①②⑤⑥⑦式可得 $F = \frac{2B^{2}d^{2}bv}{\rho}$ ⑨

当铝条的宽度 $b' > b$ 时，磁铁以速度 $v$ 进入铝条间时，磁铁受到的作用力变为 $F'$，有 $F' = \frac{2B^{2}d^{2}b'v}{\rho}$

可见，$F' > F = mg\sin\theta$，磁铁所受到的合力方向沿斜面向上，获得与运动方向相反的加速度，磁铁将减速下滑，此时加速度最大。之后，随着运动速度减小，$F'$ 也随着减小，磁铁所受的合力也减小，由于磁铁加速度与所受到的合力成正比，磁铁的加速度逐渐减小。综上所述，磁铁做加速度逐渐减小的减速运动，直到 $F' = mg\sin\theta$ 时，磁铁重新达到平衡状态，将再次以较小的速度匀速下滑。
}

\end{enumerate}

\end{document}
